\subsection{Macaulay modules and finite algebras}

Remark.--- Let \(\rho: A \to B\) be a ring homomorphism, and let \(\mathfrak{p} \in \operatorname{Spec}(A)\) . Let \(\overline{B}\) be the ring \(\kappa(\mathfrak{p}) \otimes_{A} B\) . It is identified with \(S^{-1}B/\mathfrak{p}(S^{-1}B)\) , where \(S\) is the multiplicative subset \(\rho(A - \mathfrak{p})\) of \(B\) ; the prime ideals of \(\overline{B}\) are therefore the ideals \(\mathfrak{q}\overline{B}\) , where \(\mathfrak{q}\) is a prime ideal of \(B\) which contains \(\mathfrak{p}B\) and does not meet \(S\) , in other words a prime ideal of \(B\) above \(\mathfrak{p}\) . For such an ideal \(\mathfrak{q}\) , one has \(S \subset B - \mathfrak{q}\) , hence the local ring of \(\overline{B}\) at \(\mathfrak{q}\overline{B}\) is identified with \(B_{\mathfrak{q}}/pB_{\mathfrak{q}}\) , that is, again with \(\kappa(\mathfrak{p}) \otimes_{A} B_{\mathfrak{q}}\) .

Similarly, if N is a B-module, the \(\overline{\mathrm{B}}_{\mathfrak{q}\overline{\mathrm{B}}}\) -module \(\left(\kappa(\mathfrak{p})\otimes_{\mathrm{A}}\mathrm{N}\right)_{\mathfrak{q}\overline{\mathrm{B}}}\) is identified with \(\kappa(\mathfrak{p})\otimes_{\mathrm{A}}\mathrm{N}_{\mathfrak{q}}\) . Suppose further that the B-module N is finitely generated; by Nakayama's lemma, the condition \(\kappa(\mathfrak{p})\otimes_{\mathrm{A}}\mathrm{N}_{\mathfrak{q}}=0\) is equivalent to \(\mathrm{N}_{\mathfrak{q}}=0\) . Thus the support of the \(\overline{\mathrm{B}}\) -module \(\kappa(\mathfrak{p})\otimes_{\mathrm{A}}\mathrm{N}\) is formed by the ideals \(\mathfrak{q}\overline{\mathrm{B}}\) , where \(\mathfrak{q}\) runs through the prime ideals of \(\operatorname{Supp}_{\mathrm{B}}(\mathrm{N})\) above \(\mathfrak{p}\) . In particular, for the module \(\kappa(\mathfrak{p})\otimes_{\mathrm{A}}\mathrm{N}\) to be nonzero, it is necessary and sufficient that there exists a prime ideal of \(\operatorname{Supp}_{\mathrm{B}}(\mathrm{N})\) above \(\mathfrak{p}\) .

PROPOSITION 8.--- Let \(\rho: A \to B\) be a homomorphism of Noetherian rings and let \(N\) be a \(B\) -module that is an \(A\) -module of finite type. For the \(A\) -module \(N\) to be Macaulay, it is necessary and sufficient that the \(B\) -module \(N\) be Macaulay and that, for every pair \((\mathfrak{n}, \mathfrak{n}')\) of maximal ideals of \(\mathrm{Supp}_{\mathrm{B}}(N)\) such that \(\rho^{-1}(\mathfrak{n}) = \rho^{-1}(\mathfrak{n}')\) , one has \(\dim_{B_{\mathfrak{n}}} (N_{\mathfrak{n}}) = \dim_{B_{\mathfrak{n}}'} (N_{\mathfrak{n}}')\) .

The A-module B/Ann \(_B(N)\) is isomorphic to a submodule of the finitely generated A-module \(\text{End}_A(N)\) and therefore is of finite type. Replacing A by A/Ann \(_A(N)\) and B by B/Ann \(_B(N)\) , we reduce to the case where \(\rho\) is injective and makes B a finite A-algebra, and where one has \(\text{Supp}_A(N) = \text{Spec}(A)\) and \(\text{Supp}_B(N) = \text{Spec}(B)\) . The map \(f: \text{Spec}(B) \to \text{Spec}(A)\) deduced from \(\rho\) is then surjective and a prime ideal \(\mathfrak{q}\) of B is maximal if and only if \(f(\mathfrak{q})\) is a maximal ideal of A (V, § 2, no. 1, th. 1 and prop. 1).

Let m be a maximal ideal of A. By the remark above, the prime ideals of the ring \(B_{m}\) containing \(mB_{m}\) are the ideals of the form \(qB_{m}\) where \(q \in \text{Spec}(B)\) is an ideal of B (necessarily maximal) such that \(f(\mathfrak{q}) = \mathfrak{m}\) . One has

\[
\operatorname{prof} _ {\mathrm{A} _ {\mathfrak {m}}} \left(\mathrm{N} _ {\mathfrak {m}}\right) = \operatorname{prof} _ {\mathrm{B} _ {\mathfrak {m}}} \left(\mathfrak {m B} _ {\mathfrak {m}}; \mathrm{N} _ {\mathfrak {m}}\right) = \inf _ {\mathfrak {q} \in f ^ {- 1} (\mathfrak {m})} \left(\operatorname{prof} _ {\mathrm{B} _ {\mathfrak {q}}} \left(\mathrm{N} _ {\mathfrak {q}}\right)\right)
\]

(§ 1, no. 3, prop. 4 and no. 5, prop. 8), and

\[
\dim_ {A _ {\mathfrak {m}}} \left(N _ {\mathfrak {m}}\right) = \dim_ {B _ {\mathfrak {m}}} \left(N _ {\mathfrak {m}}\right) = \sup _ {\mathfrak {q} \in f ^ {- 1} (\mathfrak {m})} \left(\dim_ {B _ {\mathfrak {q}}} \left(N _ {\mathfrak {q}}\right)\right)
\]

(VIII, § 2, no. 3, th. 1 and § 1, no. 4, prop. 9). Since we have \(\mathrm{prof}_{\mathcal{B}_{\mathfrak{q}}}(\mathbf{N}_{\mathfrak{q}}) \leqslant \dim_{\mathcal{B}_{\mathfrak{q}}}(\mathbf{N}_{\mathfrak{q}})\) for every \(\mathfrak{q} \in f^{-1}(\mathfrak{m})\) , the proposition follows from these equalities.

COROLLARY 1.--- Let \(\rho: A \to B\) be a homomorphism of Noetherian rings. If B is a finite A-algebra and a Macaulay A-module, then it is a Macaulay ring. If moreover \(\rho\) is injective, we have ht(aB) = ht(a) for every ideal a of A, and ht(b) = ht(ρ \(^{-1}\) (b)) for every ideal b of B.

The first assertion follows from prop. 8. Suppose \(\rho\) injective. Let \(\mathfrak{a}\) be an ideal of A. We have \(\mathrm{ht}(\mathfrak{a}) = \mathrm{prof}_{\mathrm{A}}(\mathfrak{a};\mathrm{B})\) since the A-module B is Macaulay, with support equal to \(\mathrm{Spec}(\mathrm{A})\) (\(n^{\circ}1\), cor. of prop. 1), \(\mathrm{ht}(\mathfrak{a}\mathrm{B}) = \mathrm{prof}_{\mathrm{B}}(\mathfrak{a}\mathrm{B};\mathrm{B})\) (loc. cit.) and \(\mathrm{prof}_{\mathrm{A}}(\mathfrak{a};\mathrm{B}) = \mathrm{prof}_{\mathrm{B}}(\mathfrak{a}\mathrm{B};\mathrm{B})\) (\(\S 1, n^{\circ}3\), prop. 4), whence \(\mathrm{ht}(\mathfrak{a}\mathrm{B}) = \mathrm{ht}(\mathfrak{a})\) . Let \(\mathfrak{b}\) be an ideal of B. By what precedes, we have \(\mathrm{ht}(\boldsymbol{\rho}^{-1}(\mathfrak{b})) = \mathrm{ht}(\boldsymbol{\rho}^{-1}(\mathfrak{b})\mathrm{B})\) . But \(\boldsymbol{\rho}^{-1}(\mathfrak{b})\mathrm{B}\) is contained in \(\mathfrak{b}\) , hence of height less than \(\mathrm{ht}(\mathfrak{b})\) and we have \(\mathrm{ht}(\mathfrak{b}) \leqslant \mathrm{ht}(\boldsymbol{\rho}^{-1}(\mathfrak{b}))\) by VIII, \(\S 2, n^{\circ}3\) , th. 1, b).

COROLLARY 2.--- Let A be an integrally closed Noetherian ring and B a ring containing A. Suppose that B is a torsion-free A-module of finite type. If B is a Macaulay ring, the A-module B is Macaulay.

Indeed, two prime ideals of B lying over the same ideal of A have the same height (VIII, § 2, no. 3, th. 2). One can therefore apply prop. 8 with N = B.

COROLLARY 3.--- Let A be an integrally closed ring, K its field of fractions, L a finite K-algebra such that {[}L : K{]} \(1_{\mathrm{A}}\) is invertible in A, and B a sub-A-algebra of L, finite over A.

\begin{enumerate}
\def\labelenumi{\alph{enumi})}
\item
  The sub-A-module A1B of B is a direct factor.
\item
  For every ideal J of A, one has the inequality \(\operatorname{prof}_{\mathrm{A}}(\mathrm{J};\mathrm{A})\geqslant\operatorname{prof}_{\mathrm{B}}(\mathrm{JB};\mathrm{B})\) .
\item
  If B is a Macaulay ring, so is A.
\end{enumerate}

The K-linear map \(\mathrm{Tr}_{\mathrm{L/K}}: \mathrm{L} \to \mathrm{K}\) maps B into A (V, § 1, no. 6, cor. 2 of prop. 17), hence defines by restriction an A-linear map \(t: \mathrm{B} \to \mathrm{A}\) . For every \(x \in \mathrm{A}\) , one has \(t(x1_{\mathrm{B}}) = [\mathrm{L}: \mathrm{K}]x\) , whence a).

By prop. 4 of § 1, no. 3, one has \(\operatorname{prof}_{\mathrm{A}}(\mathrm{J};\mathrm{B}) = \operatorname{prof}_{\mathrm{B}}(\mathrm{JB};\mathrm{B})\) ; but by a) and remark 4 of § 1, no. 1, one has \(\operatorname{prof}_{\mathrm{A}}(\mathrm{J};\mathrm{A}) \geqslant \operatorname{prof}_{\mathrm{A}}(\mathrm{J};\mathrm{B})\) , whence b).

If the ring B is Noetherian, so is A: indeed, by a) one has \(\mathfrak{a}\mathrm{B} \cap \mathrm{A} = \mathfrak{a}\) for every ideal \(\mathfrak{a}\) of A; thus every increasing sequence \((\mathfrak{a}_n)_{n \in \mathbf{N}}\) of ideals of A is stationary since the sequence \((\mathfrak{a}_n\mathrm{B})_{n \in \mathbf{N}}\) is stationary. Under the hypotheses of c), the A-module B is Macaulay (cor. 2), and the same holds for the A-module A (no. 1, example 2).

Example.--- Corollary 3 applies in particular in the following two situations:

\begin{enumerate}
\def\labelenumi{\alph{enumi})}
\item
  Consider a Noetherian integrally closed ring A, a separable extension L of its field of fractions, of finite degree n such that \(n1_{A}\) is invertible in A, and we take for B the integral closure of A in L (V, § 1, no. 6, cor. 1 of prop. 18).
\item
  Consider a Noetherian integrally closed ring B and a finite group G of automorphisms of B, such that \(\mathrm{Card(G)1_B}\) is invertible in B. We take for A the ring of elements of B invariant under the action of G. Let us verify that we are in a particular case of a). The group G acts on the field of fractions L of B, and the field of invariants of L for this action is the field of fractions K of A (V, § 1, no. 9, cor. of prop. 23). The extension L of K is Galois, and a fortiori separable; its Galois group is isomorphic to G, so that {[}L : K{]} is equal to Card G. The inverse of {[}L : K{]} \(_{1_{B}}\) is invariant under G, so that {[}L : K{]} \(_{1_{A}}\) is invertible in A. Since B is integrally closed, the ring A, equal to K ∩ B, is integrally closed and B is its integral closure in L (loc. cit., prop. 22).
\end{enumerate}

In particular, if B is a Macaulay ring, then so is A.

